\section{Implementation}
\label{sec:impl}

We implement \sys{} as an open-source Python framework (5,000+ LOC) comprising a policy engine, provenance tracker, enforcement proxy, benchmark suite, and evaluation framework.

\paragraph{System Architecture.}
The \textbf{policy engine} parses YAML policies (validated via Pydantic schemas) into compiled in-memory decision trees at startup, enabling O(1) rule lookup by tool name with short-circuit condition evaluation. Each policy rule is pre-compiled into a condition evaluator that checks provenance trust, scope constraints, temporal conditions, and rate limits in sequence, short-circuiting on the first failing condition. The engine supports hot-reloading: policy files are watched via filesystem events, and updates are applied atomically without restarting the proxy.

The \textbf{provenance tracker} maintains the information-flow DAG as an in-memory graph with PROV-DM--annotated nodes and edges (Section~\ref{sec:policies}). Each \texttt{ProvenanceNode} stores: a unique identifier, PROV-DM type (\texttt{ProvDMType}: INPUT, TOOL\_RESULT, DERIVED\_FACT, TOOL\_CALL corresponding to \texttt{prov:Entity}/\texttt{prov:Activity}), trust label from the two-element lattice $\mathcal{L} = \{\top, \bot\}$, SHA-256 content hash, timestamp, parent edges annotated with PROV-DM relations (\texttt{wasDerivedFrom}, \texttt{wasGeneratedBy}, \texttt{used}), and a 384-dimensional sentence embedding (all-MiniLM-L6-v2, computed lazily on first provenance query). Trust propagation uses an explicit \texttt{trust\_meet(*labels)} operator implementing the lattice meet $\sqcap$: any $\bot$ ancestor \emph{or} unresolved provenance yields $\bot$ (conservative default). Crucially, LLM-generated content with \emph{no traceable source nodes} also defaults to $\bot$---the conservative choice that a value with no provenance should not be assumed safe. Only values with at least one explicit trusted ancestor and no untrusted ancestors receive label $\top$. Provenance resolution follows the three-stage pipeline (Section~\ref{sec:policies}): substring match $\to$ embedding cosine similarity $\to$ conservative default ($\bot$). The embedding model adds $\sim$90\,MB memory and $<$5\,ms per query; it can be disabled via \texttt{PROVSAFE\_EMBEDDING=0} for resource-constrained deployments (falling back to substring + conservative default). An explicit \texttt{declassify()} method implements controlled trust upgrade ($\bot \to \top$) exclusively through authenticated user confirmation, enforcing the Declassification Control property.

The \textbf{enforcement proxy} implements the seven-stage workflow (Section~\ref{sec:proxy}) with three defense layers: pre-LLM input screening (47-rule regex scanner for jailbreaks, encoding detection, chaining patterns), argument validation (path traversal, command injection, encoded payload detection), and policy evaluation. Critically, the argument validation stage \emph{marks suspicious arguments as UNTRUSTED} rather than short-circuiting to an immediate denial---validated arguments continue through the provenance and policy stages, which make the final allow/deny/confirm decision. This pipeline integration ensures that every defense layer is consulted on every tool call and that the contribution of each layer is isolable in ablation studies. Decision logging uses append-only JSON with SHA-256 hash chaining for tamper-evident auditability. The proxy exposes both a Python decorator interface (for in-process agent frameworks) and an HTTP/gRPC server interface (for remote agents), requiring no modifications to agent internals.

\paragraph{Benchmark Suite.}
The benchmark contains 160 attack scenarios across 8 categories (Direct Injection: 20, Device-State Injection: 22, File Content Injection: 22, Multi-Turn Chaining: 20, Encoding Obfuscation: 18, Role-Play Jailbreak: 20, Confused Deputy: 18, Privilege Escalation: 20) and 40 benign tasks spanning device control, calendar management, file operations, sensor queries, and notification dispatch. Each scenario specifies: an attack description, the injection vector, expected malicious tool call, ground-truth label (malicious/benign), and the attack category. All 11 mock tools (file system, calendar, notification, email, device control) use deterministic in-memory backends with seeded RNGs (seed 42) for reproducibility. The evaluation framework orchestrates comparative evaluation of four configurations (No Defense, Pattern Filter, Policy-Only, \sys{} Full) and ablation studies, producing structured JSON results with per-scenario outcomes.

\paragraph{Policy Configurations.}
Three policies ship with \sys{}:
\begin{itemize}[leftmargin=1.2em]
  \item \emph{Default}: The policy used in all evaluations. Defines risk tiers (LOW: read-only, MEDIUM: create/modify, HIGH: destructive/exfiltrating, CRITICAL: system-level), provenance restrictions (untrusted provenance requires confirmation for HIGH/CRITICAL tools), scope constraints, and dangerous-action hardcoded blocks.
  \item \emph{Permissive}: All tool calls are allowed regardless of provenance. Equivalent to the No Defense baseline.
  \item \emph{Strict}: All tool calls with any untrusted provenance are denied. Achieves near-0\% ASR-IA but severely restricts usability, demonstrating the cost of binary taint semantics without provenance-gated policies.
\end{itemize}
Operators customize via YAML; tools register declaratively with name, risk tier, typed schema, and optional scope constraints.

\paragraph{Performance.}
Proxy overhead is sub-millisecond: policy evaluation median 0.06\,ms (P95: 0.12\,ms), provenance queries median 0.82\,ms (P95: 1.8\,ms with 500+ graph nodes), argument validation median 0.15\,ms (P95: 0.3\,ms). End-to-end task runtime averages 20.21\,s vs.\ 12.35\,s undefended; the 7.86\,s difference is dominated by user confirmation delays (average 5.2\,s per confirmation) and LLM re-prompting after blocked calls (average 2.1\,s per re-prompt), not proxy computation. Memory footprint is $\sim$55\,MB with 1,000 provenance nodes; the embedding extension adds $\sim$90\,MB for the all-MiniLM-L6-v2 model weights.

\paragraph{Evaluation Infrastructure.}
The primary experiments use four local models served via LM Studio v0.3+ on an Apple M4 Pro workstation (24\,GB unified RAM) with OpenAI-compatible REST API (\texttt{localhost:1234}) and automatic model load/unload: Meta-Llama-3.1-8B-Instruct (8B, Q4\_K\_M, 4.92\,GB), Qwen2.5-7B-Instruct (7B, Q4\_K\_M, 4.68\,GB), Gemma-2-9B-IT (9B, Q4\_K\_M, 5.76\,GB), and Phi-3.5-Mini-Instruct (3B, Q4\_K\_M, 2.18\,GB). Large-model validation uses Qwen3.5-122B served via Open WebUI on a university GPU cluster. Repetition~1 uses temperature $T=0.0$ (greedy decoding); repetitions~2--5 use $T \in \{0.05, 0.10, 0.15, 0.20\}$ with unique per-trial seeds derived from a cryptographic hash of (scenario, system, model, repetition), ensuring statistically independent trials for bootstrap confidence interval validity. The embedding model (all-MiniLM-L6-v2) runs on CPU with $<$5\,ms per embedding. Over 5 repetitions of 200 scenarios $\times$ 5 models $\times$ 5 systems, this yields 25{,}000 total trials.

\paragraph{Availability.}
Code will be released as open-source with documentation, policies, the full benchmark suite, and reproducible evaluation scripts. Installation: \texttt{pip install provsafe}.
